\vspace{0.5em}\noindent\textbf{From Static Recommendations to Real-time Personalization: A
Journey with Amazon
Personalize}\par

When first building the ``Recommended Products'' section for an
e-commerce site, I thought it was simple: just take the Top Selling or
Trending items of the week and display them to all users. But as I
tracked conversion metrics, I realized a truth: static recommendations
do not understand the individual behavior of each customer. A user who
just viewed running shoes would still be recommended\ldots{} the
best-selling rice cooker of the week.

The challenge was: how do we transition from ``one-size-fits-all
recommendations'' to real-time personalized recommendations, when the
team doesn't have an in-house Data Science crew to train Collaborative
Filtering or Deep Learning models? The answer I chose was \textbf{Amazon
Personalize} - a managed ML service from AWS that uses the exact same
technology behind Amazon.com's recommendation system.

Below is what I learned after implementing it in production.

\vspace{0.5em}\noindent\textbf{1. Data Preparation: The Three Pillars - Interactions, Users,
Items}\par

\textbf{Reality:} Many teams rush into creating a Dataset Group,
forgetting that input data quality determines 80\% of the model's
quality. Amazon Personalize requires data across three types of
datasets, but only one is mandatory: - \textbf{Item interactions
(mandatory):} records user behaviors - clicks, add-to-carts, purchases.
Requires at least three fields: \texttt{USER\_ID}, \texttt{ITEM\_ID},
\texttt{TIMESTAMP}, and should include \texttt{EVENT\_TYPE} to
distinguish between different types of behaviors. - \textbf{Users
(optional):} only requires the \texttt{USER\_ID} field, can include up
to 5 metadata fields such as age, region, gender. - \textbf{Items
(optional):} only requires \texttt{ITEM\_ID}, can include metadata like
category, price, brand - this helps the model handle new products better
(cold-start).

\textbf{Lesson:} Regarding the minimum data threshold for meaningful
training, AWS recommends having at least 1,000 users and 1,000
interactions, and significantly more data per product to achieve better
results. For the User-Personalization recipe, official AWS documentation
also explicitly states a minimum requirement of 1,000 interaction
records.

Another important point: if your Interactions dataset has multiple event
types (click, purchase, watch\ldots), you can specify a particular event
type to train on, or - with v2 recipes - assign different weights to
each event type (e.g., prioritizing purchases over clicks). This is a
detail I overlooked initially, causing the model to ``learn'' too much
from random click behaviors instead of truly valuable purchasing
behaviors.

\vspace{0.5em}\noindent\textbf{2. Dataset Group, Schema \& Import: The Foundation Before
Training}\par

\textbf{Reality:} Amazon Personalize organizes everything around the
concept of a Dataset Group - a container holding a maximum of three
dataset types (Interactions, Items, Users, and for more advanced use
cases, Actions/Action interactions). You can choose between: -
\textbf{Domain dataset group} (ECOMMERCE, VIDEO\_ON\_DEMAND\ldots): uses
pre-configured recommenders, suitable for rapid deployment. -
\textbf{Custom dataset group:} lets you choose recipes and configure
Solutions yourself --- more flexible for specific use cases.

\textbf{Lesson:} Define the schema (Avro-based) for each dataset clearly
from the start. If the team is familiar with IaC, you can declare
everything using AWS CDK instead of clicking manually on the Console ---
this allows for easily reproducing Staging/Production environments and
auditing when issues occur. A practical CDK structure example:

\begin{Shaded}
\begin{Highlighting}[]
\NormalTok{DatasetGroup (domain: ECOMMERCE)}
\NormalTok{ └── InteractionsDataset (schema: USER\_ID, ITEM\_ID, TIMESTAMP, EVENT\_TYPE)}
\NormalTok{ └── ItemsDataset (schema: ITEM\_ID, CATEGORY, PRICE...)}
\NormalTok{ └── UsersDataset (schema: USER\_ID, AGE, REGION...)}
\end{Highlighting}
\end{Shaded}

After creating the datasets, data import can be done via Batch import
(from S3, suitable for historical data) or Streaming ingestion via
PutEvents API (suitable for real-time behaviors like clicks on the
site).

\vspace{0.5em}\noindent\textbf{3. Solution \& Recipe: Choosing the Right ``Brain'' for the
Problem}\par

\textbf{Reality:} This is the step where newcomers get most confused ---
Amazon Personalize provides multiple recipes (pre-packaged algorithms)
divided into three groups: - \textbf{USER\_PERSONALIZATION} --- predicts
the products a user is most likely to interact with next (used for
homepages, recommendation emails). - \textbf{PERSONALIZED\_RANKING} ---
re-ranks an existing list of products based on relevance to each user
(used for category/search pages). - \textbf{RELATED\_ITEMS} ---
recommends ``similar products'' (used for product detail pages).

\textbf{Lesson:} AWS currently recommends using v2 recipes (e.g.,
\texttt{User-Personalization-v2}) instead of the original versions, as
these recipes can handle up to 5 million products with faster training
times and lower latency. In addition, the User-Personalization recipe
features an Exploration mechanism - proactively inserting new or
sparsely interacted products into recommendations to avoid the ``filter
bubble'', which is highly useful when the catalog changes rapidly (new
products arriving constantly).

If you need a ``no model training required'' option for scenarios like
hourly top trending items, Amazon Personalize also offers the
\texttt{Trending-Now} recipe. This helps identify products with the
fastest-growing interactions over a recent period - a great supplement
for a ``Trending Now'' section alongside the main personalized
recommendations.

\vspace{0.5em}\noindent\textbf{4. Campaign: Serving the Model in
Real-time}\par

\textbf{Reality:} Having a solution version (a trained model) is not
enough - you need a Campaign to create an actual endpoint that returns
real-time results when the application calls the
\texttt{GetRecommendations} or \texttt{GetPersonalizedRanking} API.

\textbf{Most important lesson in this step:} real-time personalization
doesn't stop at deploying the Campaign. With the User-Personalization
recipe, Amazon Personalize automatically updates the latest model every
two hours to consider new products - meaning you don't have to write a
manual retraining pipeline just to add new products to the catalog.
Furthermore, when there are filters (e.g., removing out-of-stock
products from recommendations), these filters are also updated within 15
minutes of the latest data import or incremental record, fast enough for
most e-commerce needs.

\vspace{0.5em}\noindent\textbf{Overall Architecture I
Applied}\par

\begin{Shaded}
\begin{Highlighting}[]
\NormalTok{Clickstream / Purchase Events (Frontend, Backend)}
\NormalTok{ │ (PutEvents API {-} real{-}time)}
\NormalTok{ ▼}
\NormalTok{Amazon Personalize — Interactions Dataset}
\NormalTok{ │}
\NormalTok{ ├── Dataset Group (ECOMMERCE domain)}
\NormalTok{ ├── Solution (recipe: User{-}Personalization{-}v2)}
\NormalTok{ └── Campaign (real{-}time endpoint)}
\NormalTok{ │}
\NormalTok{ ▼}
\NormalTok{GetRecommendations API}
\NormalTok{ │}
\NormalTok{ ▼}
\NormalTok{Homepage / Product Pages / Email marketing}
\end{Highlighting}
\end{Shaded}

\textbf{Result:} After switching from static Top Selling to personalized
recommendations, the ``You might also like'' sections on the homepage
and product detail pages display correctly according to each user's
behavior. We no longer need a Data Science team on standby to fine-tune
the model manually - the entire train/deploy/update lifecycle is managed
by Amazon Personalize.

\vspace{0.5em}\noindent\textbf{Biggest Takeaway}\par

The tool is only half the story. The other half lies in preparing
behavioral data correctly and choosing the right recipe for the right
display location (the homepage is different from a category page, which
is different from a product detail page). A good personalized
recommendation system isn't the most complex model, but a model fed with
clean interaction data, continuously updated, and placed in the right
spot in the customer journey.

\vspace{0.5em}\noindent\textbf{References (Official AWS
Docs)}\par

For those who want to dive deeper into technical details:

\vspace{0.5em}\noindent\textbf{Dataset \& Schema:}\par

\begin{itemize}
\tightlist
\item
  \href{https://docs.aws.amazon.com/personalize/latest/dg/event-values-types.html}{Amazon
  Personalize Developer Guide --- Choosing item interaction data for
  training}
\item
  \href{https://docs.aws.amazon.com/personalize/latest/dg/API_CreateDatasetGroup.html}{CreateDatasetGroup
  API Reference}
\end{itemize}

\vspace{0.5em}\noindent\textbf{Recipe \& Solution:}\par

\begin{itemize}
\tightlist
\item
  \href{https://docs.aws.amazon.com/personalize/latest/dg/native-recipe-new-item-USER_PERSONALIZATION.html}{User-Personalization
  Recipe Guide}
\item
  \href{https://aws.amazon.com/blogs/machine-learning/recommend-top-trending-items-to-your-users-using-the-new-amazon-personalize-recipe/}{Trending-Now
  Recipe (AWS ML Blog)}
\end{itemize}

\vspace{0.5em}\noindent\textbf{Operations \& Data Updates:}\par

\begin{itemize}
\tightlist
\item
  \href{https://aws.amazon.com/blogs/machine-learning/simplify-data-management-with-new-apis-in-amazon-personalize}{Simplify
  Data Management with New APIs in Amazon Personalize (AWS ML Blog)}
\end{itemize}

\textbf{Original post}:
\href{https://www.facebook.com/groups/awsstudygroupfcj/permalink/2229289067836053/?rdid=PGigWM6loaBLr4ZG\#}{Facebook
AWS Study Group FCJ}
